% Part: sets-functions-relations
% Chapter: arithmetization
% Section: reflections

\documentclass[../../../include/open-logic-section]{subfiles}

\begin{document}

\olfileid[fr]{sfr}{arith}{ref}
\olsection{Quelques réflexions philosophiques}

Voilà pour les aspects techniques. Mais qu'avons-nous obtenu ?

D'abord, et cela ne prête guère à contestation, de très belles
mathématiques pures. Nous avons aussi obtenu des résultats conceptuels
profonds. Comprendre que la propriété de la borne supérieure exprime
la différence cruciale entre réels et rationnels constitue une avancée
majeure. De plus, la construction explicite des réels comme coupures de
Dedekind donne à l'analyse un fondement solide. Nous savons que la
notion de \emph{corps ordonné complet} est cohérente, puisque les
coupures forment précisément un tel corps.
\begin{editorial}
Cette dernière affirmation est relative au cadre ensembliste dans
lequel la construction est menée : celle-ci réalise les propriétés
d'un corps ordonné complet. Elle ne démontre pas, à elle seule,
la cohérence absolue de la théorie des ensembles employée.
\end{editorial}

Ces résultats appellent néanmoins quelques réserves.

Premièrement, il n'est pas clair qu'envisager les réels comme des
coupures soit \emph{plus} rigoureux que de les envisager à partir de
leurs développements décimaux familiers, éventuellement infinis.
Cette dernière « construction » ressemble quelque peu à celle par
les suites de Cauchy ; pour l'essentiel, elle était en fait connue
des mathématiciens dès le début du dix-septième siècle
(voir la \olref[cauchy]{sec}). Le véritable gain de rigueur est venu de
la reconnaissance de la propriété de la borne supérieure.
La possibilité de construire les réels comme des ensembles particuliers
n'est peut-être pas, à elle seule, si intéressante.

Il est encore moins clair que l'arithmétisation, beaucoup plus facile,
des entiers relatifs ou des rationnels accroisse la rigueur dans ces
domaines. Une observation simple mérite ici d'être faite. Après avoir
\emph{construit} les entiers relatifs comme classes d'équivalence de
couples ordonnés de naturels, puis les rationnels comme classes
d'équivalence de couples ordonnés d'entiers relatifs, puis les réels
comme ensembles de rationnels, nous \emph{oublions} immédiatement
ces constructions. En particulier, personne ne voudrait les
\emph{invoquer} dans une démonstration mathématique, sauf, bien sûr,
pour démontrer qu'elles ont les propriétés attendues. Il est bien
plus facile de parler directement d'un réel que d'un ensemble
d'ensembles d'ensembles d'ensembles d'ensembles d'ensembles d'ensembles
de naturels.
% (Les réels ont été construits comme ensembles de rationnels ;
% ceux-ci ont été construits comme classes d'équivalence de couples
% ordonnés d'entiers relatifs, c'est-à-dire comme ensembles d'ensembles
% d'ensembles d'entiers relatifs ; et ainsi de suite.)
% 2/3 = [2,3]
% c'est-à-dire {<2,3>, ... }
% { {{2},{2,3}}, ... }
%
% les réels sont
% des ensembles de rationnels
%
% les rationnels sont
% des ensembles d'ensembles d'ensembles d'entiers relatifs
%
% les entiers relatifs sont
% des ensembles d'ensembles d'ensembles de naturels

Plus douteux encore est que ces définitions nous disent ce que
\emph{sont} les entiers relatifs, les rationnels ou les réels,
\emph{sur le plan métaphysique}. Autrement dit, il est douteux que
les réels, par exemple, \emph{soient} certains ensembles
(d'ensembles d'ensembles\ldots). Le principal obstacle est que
la construction aurait pu être menée de bien des façons. Dans le
cas des réels, il existe des constructions différentes qui présentent
un réel intérêt (voir la \olref[cauchy]{sec}). Mais voici un moyen
tout à fait élémentaire d'obtenir d'autres constructions : comme
dans la \olref[sfr][rel][ref]{sec}, nous aurions pu définir les couples
ordonnés un peu autrement. Avec cette autre définition, nos
constructions auraient tout aussi bien fonctionné, mais nous aurions
obtenu d'autres objets. Il existe donc de nombreuses constructions
ensemblistes concurrentes des entiers relatifs, des rationnels et
des réels. Affirmer que les entiers relatifs, par exemple, sont
\emph{ces} ensembles plutôt que \emph{ceux-là} serait alors arbitraire,
et embarrassant. Comme dans la \olref[sfr][rel][ref]{sec}, il s'agit
d'un cas de l'argument rendu célèbre par \citealt{Benacerraf1965}.

Un autre point mérite d'être soulevé : nos constructions ont quelque
chose d'assez \emph{étrange}. Nous sommes partis des nombres naturels.
Nous construisons ensuite les entiers relatifs et le « $0$ des entiers
relatifs », c'est-à-dire $\equivrep{0,0}{\Intequiv}$.
Or, dans la réalisation ensembliste usuelle,
$0\neq\equivrep{0,0}{\Intequiv}$ ; avec ces constructions,
\emph{aucun} naturel n'est un entier relatif ainsi construit.
Cela paraît pourtant très contraire à l'intuition.
Dans la \olref[sfr][set][imp]{sec}, nous avons même affirmé,
sans véritable justification, que $\Nat\subseteq\Rat$.
Si ces constructions nous disent exactement \emph{ce que sont}
les nombres, cette affirmation, prise littéralement dans cette
réalisation, est manifestement fausse.
\begin{editorial}
Les inégalités et la séparation littérale invoquées ici dépendent du
codage des objets. Elles valent, par exemple, si les naturels sont les
ordinaux finis de von Neumann : les classes de couples qui représentent
les entiers relatifs et les rationnels sont infinies, tandis que chaque
naturel est un ensemble fini. La source laisse ce choix de réalisation
implicite. En revanche, les plongements déjà définis permettent
d'identifier chaque système de nombres à son image dans le suivant ;
ils justifient les inclusions usuelles après cette identification.
La distinction pertinente est celle entre égalité littérale des objets
et identification au moyen d'un plongement.
\end{editorial}

Prenons donc du recul. En travaillant dans une théorie naïve des
ensembles et en nous donnant les naturels, nous pouvons \emph{traiter}
les entiers relatifs, les rationnels et les réels comme certains
ensembles. En ce sens, nous pouvons \emph{plonger} les théories de
ces objets dans une théorie des ensembles. La portée philosophique
de ce plongement n'est cependant pas si simple à déterminer.

Bien sûr, ce n'est pas le dernier mot !{} Il s'agit seulement de
souligner ceci : montrer que l'arithmétisation des réels \emph{a}
une portée philosophique profonde exigerait un argument
\emph{philosophique} supplémentaire.

% De plus, tout au long de ce chapitre, nous avons tenu les naturels
% pour simplement donnés. Mais que pouvons-nous faire à leur sujet ?
% C'est l'objet du chapitre suivant.

\end{document}
